% =============================================================================
\section{Introduction et contexte}
% =============================================================================

% Les trois ressources du fil rouge et leur interdépendance.
\begin{frame}{Le bloc opératoire : des ressources rares et couplées}
  \small
  \begin{columns}[T,onlytextwidth]
    \begin{column}{0.31\textwidth}
      \begin{exampleblock}{\centering Bloc opératoire}
        \centering\footnotesize
        Salles, vacations de 4\,h\\
        Chirurgien, équipe d'anesthésie\\
        Temps opératoire contraint
      \end{exampleblock}
    \end{column}
    \begin{column}{0.31\textwidth}
      \begin{block}{\centering Lits d'hospitalisation}
        \centering\footnotesize
        Capacité limitée ($\sim$45 lits)\\
        Séjours de plusieurs jours\\
        Pics et creux d'occupation
      \end{block}
    \end{column}
    \begin{column}{0.31\textwidth}
      \begin{alertblock}{\centering Places ambulatoires}
        \centering\footnotesize
        Séjours courts\\
        Forte sensibilité à l'heure\\
        Ressource distincte des lits
      \end{alertblock}
    \end{column}
  \end{columns}

  \vspace{0.35cm}
  \begin{alertblock}{Le constat}
    \small Ces décisions sont \textbf{interdépendantes} : un lit non disponible
    annule une opération, une place ambulatoire saturée bloque la sortie de
    salle, une salle fermée déplace toute la journée. Pourtant, ces ressources
    sont le plus souvent \textbf{pilotées séparément}.
  \end{alertblock}
  \source{Dimensionnement clinique : $\sim$45 lits classiques et 21 lits
    ambulatoires (consignes du fil rouge).}
\end{frame}

% Le glissement de logique : flux poussé -> flux tiré.
\begin{frame}{Du flux poussé au flux tiré}
  \small
  \begin{columns}[T]
    \begin{column}{0.48\textwidth}
      \begin{block}{Organisation actuelle : flux poussé}
        \begin{itemize}
          \item Le chirurgien choisit une date selon son planning.
          \item Blocs, lits et équipes s'adaptent ensuite.
          \item Résultat : pics et creux d'activité.
        \end{itemize}
      \end{block}
    \end{column}
    \begin{column}{0.48\textwidth}
      \begin{block}{Approche proposée : flux tiré}
        \begin{itemize}
          \item La date dépend des ressources disponibles.
          \item Bloc et lits sont optimisés conjointement.
          \item Le chirurgien choisit parmi deux propositions.
        \end{itemize}
      \end{block}
    \end{column}
  \end{columns}

  \vspace{0.2cm}
  \begin{alertblock}{Idée centrale du projet}
    Optimiser \textbf{simultanément} le bloc, les lits et l'ambulatoire, et
    proposer au praticien une \textbf{Date A} (optimale) et une
    \textbf{Date B} (alternative).
  \end{alertblock}
  \source{Le chirurgien conserve la décision finale : l'outil est une aide à la
    décision, pas un automate.}
\end{frame}

